\documentclass[a4paper,12pt]{article}
\usepackage[utf8]{inputenc}
\usepackage[T2A]{fontenc}
\usepackage[russian]{babel}
\usepackage{amsmath,amssymb,amsthm}
\usepackage{graphicx}
\usepackage{booktabs}
\usepackage{hyperref}
\usepackage{listings}
\usepackage{xcolor}
\usepackage{algorithm}
\usepackage{algpseudocode}
\usepackage{geometry}
\usepackage{float}
\usepackage{caption}
\usepackage{subcaption}

\geometry{margin=2.5cm}

% Настройка листингов кода
\definecolor{codegreen}{rgb}{0,0.6,0}
\definecolor{codegray}{rgb}{0.5,0.5,0.5}
\definecolor{codepurple}{rgb}{0.58,0,0.82}
\definecolor{backcolour}{rgb}{0.95,0.95,0.92}

\lstdefinestyle{mystyle}{
    backgroundcolor=\color{backcolour},
    commentstyle=\color{codegreen},
    keywordstyle=\color{magenta},
    numberstyle=\tiny\color{codegray},
    stringstyle=\color{codepurple},
    basicstyle=\ttfamily\footnotesize,
    breakatwhitespace=false,
    breaklines=true,
    captionpos=b,
    keepspaces=true,
    numbers=left,
    numbersep=5pt,
    showspaces=false,
    showstringspaces=false,
    showtabs=false,
    tabsize=2
}
\lstset{style=mystyle}

\title{Сравнительный анализ оптимизаторов AdamW, Muon и гибридного подхода для параметрически-эффективного дообучения больших языковых моделей}
\author{Рштуни В.}
\date{\today}

\begin{document}

\maketitle

\begin{abstract}
В данном отчёте представлено детальное исследование трёх стратегий оптимизации градиентного спуска для параметрически-эффективного дообучения (LoRA) больших языковых моделей. Рассматриваются: классический адаптивный оптимизатор AdamW, современный оптимизатор Muon на основе матричной ортогонализации, и предложенный гибридный подход, комбинирующий преимущества обоих методов. Для каждого оптимизатора приведены математические формулы, алгоритмы в псевдокоде и анализ вычислительной сложности. Эксперименты проведены на модели Qwen2.5-0.5B с использованием датасета OpenWebText и оценкой на задаче логического вывода PIQA. Результаты показывают, что AdamW обеспечивает наилучшую скорость обучения (0.682 шага/с), Muon экономит $\sim$2\% видеопамяти, а гибридный подход демонстрирует сбалансированные характеристики. Все методы достигают сопоставимого качества на тестовой выборке ($\sim$70\% accuracy), что указывает на достаточность данных методов для задач малого масштаба.
\end{abstract}
\newpage

\tableofcontents
\newpage

\section{Введение}
\label{sec:intro}

Современные большие языковые модели (LLM) содержат сотни миллионов и миллиарды параметров, что делает их полное дообучение вычислительно дорогостоящим. Методы параметрически-эффективного дообучения, такие как LoRA (Low-Rank Adaptation)~\cite{hu2021lora}, позволяют адаптировать предобученные модели к новым задачам, обновляя лишь малую долю параметров (обычно $<$5\%).

Ключевым компонентом процесса обучения является выбор оптимизатора --- алгоритма, определяющего направление и величину обновления весов модели на основе градиентов функции потерь. В последние годы появились новые методы оптимизации, претендующие на улучшение сходимости и эффективности использования ресурсов по сравнению с классическими подходами.

\textbf{Цель работы}: провести сравнительный анализ трёх оптимизаторов для LoRA-дообучения:
\begin{enumerate}
    \item \textbf{AdamW}~\cite{loshchilov2019adamw} --- стандартный адаптивный метод с коррекцией веса;
    \item \textbf{Muon}~\cite{moonshot2024muon} --- оптимизатор на основе матричной ортогонализации;
    \item \textbf{Hybrid} --- авторский гибридный подход, комбинирующий Muon и AdamW.
\end{enumerate}

\textbf{Критерии сравнения}: время обучения, пиковое потребление GPU-памяти, финальное значение функции потерь, точность на тестовой задаче (PIQA), скорость сходимости.
\newpage

\section{Теоретические основы}
\label{sec:theory}

\subsection{Постановка задачи обучения}

Пусть дана модель с параметрами $\theta \in \mathbb{R}^d$ и функция потерь $\mathcal{L}(\theta; \mathcal{D})$, вычисляемая на обучающем датасете $\mathcal{D}$. Цель обучения --- найти параметры $\theta^*$, минимизирующие эмпирический риск:
\begin{equation}
    \theta^* = \arg\min_{\theta} \mathbb{E}_{(x,y) \sim \mathcal{D}}[\ell(f_\theta(x), y)],
\end{equation}
где $\ell$ --- функция потерь для отдельного примера, $f_\theta$ --- предсказание модели.

В контексте дообучения с помощью LoRA исходные веса $W_0$ замораживаются, а к ним добавляется низкоранговое приращение:
\begin{equation}
    W = W_0 + \Delta W = W_0 + BA, \quad A \in \mathbb{R}^{r \times k}, \; B \in \mathbb{R}^{n \times r}, \; r \ll \min(n,k).
\end{equation}
Обучению подлежат только матрицы $A$ и $B$, что существенно снижает количество обучаемых параметров.

\subsection{Проблемы стандартного стохастического градиентного спуска}

Классический SGD обновляет параметры по правилу:
\begin{equation}
    \theta_{t+1} = \theta_t - \eta \cdot g_t, \quad g_t = \nabla_\theta \mathcal{L}(\theta_t),
\end{equation}
где $\eta$ --- скорость обучения.

Основные ограничения SGD:
\begin{itemize}
    \item \textbf{Чувствительность к масштабу градиентов}: параметры с разными масштабами требуют разных шагов обучения;
    \item \textbf{Шум мини-батчей}: стохастичность градиентов может приводить к колебаниям траектории;
    \item \textbf{Отсутствие адаптации}: единая скорость обучения для всех параметров неэффективна для разреженных или редко обновляемых весов.
\end{itemize}

Эти проблемы мотивируют использование адаптивных методов оптимизации.
\newpage

\section{Описание оптимизаторов}
\label{sec:optimizers}

\subsection{AdamW: адаптивный момент с коррекцией веса}

\subsubsection{Идея метода}

AdamW~\cite{loshchilov2019adamw} сочетает преимущества двух подходов:
\begin{itemize}
    \item \textbf{Momentum}: накопление экспоненциально затухающего среднего градиентов для сглаживания обновлений;
    \item \textbf{Adaptive scaling}: нормировка градиентов на оценку их дисперсии для автоматической адаптации шага по каждому параметру;
    \item \textbf{Weight decay decoupling}: применение регуляризации отдельно от градиентного шага, что улучшает обобщение.
\end{itemize}

\subsubsection{Математическая формулировка}

На шаге $t$ для каждого параметра $\theta_i$ вычисляются:

\begin{align}
    m_t &= \beta_1 m_{t-1} + (1 - \beta_1) g_t \quad &\text{(первый момент)} \\
    v_t &= \beta_2 v_{t-1} + (1 - \beta_2) g_t^2 \quad &\text{(второй момент)} \\
    \hat{m}_t &= \frac{m_t}{1 - \beta_1^t} \quad &\text{(коррекция смещения)} \\
    \hat{v}_t &= \frac{v_t}{1 - \beta_2^t} \quad &\text{(коррекция смещения)} \\
    \theta_{t+1} &= \theta_t - \eta \left( \frac{\hat{m}_t}{\sqrt{\hat{v}_t} + \epsilon} + \lambda \theta_t \right) \quad &\text{(обновление с weight decay)}
\end{align}

где:
\begin{itemize}
    \item $g_t$ --- градиент на шаге $t$;
    \item $\beta_1, \beta_2 \in [0, 1)$ --- коэффициенты затухания моментов (обычно 0.9 и 0.999);
    \item $\epsilon > 0$ --- малая константа для численной стабильности ($10^{-8}$);
    \item $\lambda$ --- коэффициент weight decay;
    \item $\eta$ --- базовая скорость обучения.
\end{itemize}

\subsubsection{Псевдокод}

\begin{algorithm}[H]
\caption{AdamW (упрощённая версия)}
\label{alg:adamw}
\begin{algorithmic}[1]
\Require Параметры $\theta$, градиенты $g$, гиперпараметры $\eta, \beta_1, \beta_2, \epsilon, \lambda$
\State Инициализировать $m \gets 0$, $v \gets 0$, $t \gets 0$
\While{не достигнут критерий остановки}
    \State $t \gets t + 1$
    \State $g \gets \nabla_\theta \mathcal{L}(\theta)$ \Comment{Вычисление градиента}
    \State $m \gets \beta_1 \cdot m + (1 - \beta_1) \cdot g$
    \State $v \gets \beta_2 \cdot v + (1 - \beta_2) \cdot g^2$
    \State $\hat{m} \gets m / (1 - \beta_1^t)$
    \State $\hat{v} \gets v / (1 - \beta_2^t)$
    \State $\theta \gets \theta - \eta \cdot \left( \frac{\hat{m}}{\sqrt{\hat{v}} + \epsilon} + \lambda \cdot \theta \right)$
\EndWhile
\State \Return $\theta$
\end{algorithmic}
\end{algorithm}

\subsubsection{Вычислительная сложность и память}

\begin{itemize}
    \item \textbf{Память}: $2d$ дополнительных скаляров на параметр ($m$ и $v$);
    \item \textbf{Вычисления}: $O(d)$ операций на шаг (элементарные покомпонентные операции);
    \item \textbf{Параллелизация}: тривиальная, все операции покомпонентные.
\end{itemize}

\subsection{Muon: оптимизация через матричную ортогонализацию}

\subsubsection{Идея метода}

Muon (Matrix Orthogonalization for Updating Networks)~\cite{moonshot2024muon} основан на гипотезе, что направление градиента можно улучшить путём ортогонализации матрицы обновлений. Метод применяет итерации Ньютона-Шульца для приближения полярного разложения градиентной матрицы, что теоретически ускоряет сходимость в задачах с коррелированными градиентами.

Ключевые особенности реализации в данном проекте:
\begin{itemize}
    \item Применяется только к \textbf{двумерным тензорам} (весовым матрицам), исключая эмбеддинги и \texttt{lm\_head};
    \item Использует \textbf{один буфер момента} вместо двух, экономя память;
    \item Ортогонализация выполняется через \textbf{квинтичные итерации Ньютона-Шульца} с коэффициентами $(a, b, c) = (3.4445, -4.7750, 2.0315)$;
    \item Автоматическая корректировка learning rate в зависимости от размера матрицы: $\eta_{\text{adj}} = \eta \cdot 0.2 \cdot \sqrt{\max(n, m)}$.
\end{itemize}

\subsubsection{Математическая формулировка}

Пусть $G \in \mathbb{R}^{n \times m}$ --- матрица градиентов для некоторого весового слоя.

\textbf{Шаг 1: Нормировка и подготовка}
\begin{equation}
    X_0 = \frac{G}{\|G\|_F + \epsilon}, \quad \text{если } n > m \text{, то } X_0 \gets X_0^\top
\end{equation}

\textbf{Шаг 2: Квинтичные итерации Ньютона-Шульца}

Для $k = 0, 1, \ldots, K-1$:
\begin{align}
    A_k &= X_k X_k^\top \\
    B_k &= b \cdot A_k + c \cdot A_k^2 \\
    X_{k+1} &= a \cdot X_k + B_k X_k
\end{align}
где коэффициенты $(a, b, c) = (3.4445, -4.7750, 2.0315)$ подобраны для максимизации скорости сходимости.

\textbf{Шаг 3: Обновление с моментом и корректировкой шага}
\begin{align}
    M_t &= \beta M_{t-1} + (1 - \beta) \cdot X_K \\
    \eta_{\text{adj}} &= \eta \cdot 0.2 \cdot \sqrt{\max(n, m)} \\
    \theta_{t+1} &= \theta_t - \eta_{\text{adj}} \cdot M_t - \eta \cdot \lambda \cdot \theta_t
\end{align}

\subsubsection{Псевдокод}

\begin{algorithm}[H]
\caption{Muon: обновление одного весового матричного параметра}
\label{alg:muon}
\begin{algorithmic}[1]
\Require Матрица весов $W$, градиент $G$, гиперпараметры $\eta, \beta, \lambda, K, (a,b,c)$
\If{$G$ не является 2D-тензором}
    \State \Return AdamW-обновление для $G$ \Comment{Fallback для векторов}
\EndIf
\State \Comment{Шаг 1: нормировка}
\State $X \gets G / (\|G\|_F + \epsilon)$
\If{$\text{rows}(G) > \text{cols}(G)$}
    \State $X \gets X^\top$
\EndIf
\State \Comment{Шаг 2: квинтичные итерации Ньютона-Шульца}
\For{$k = 1$ \textbf{to} $K$}
    \State $A \gets X \cdot X^\top$
    \State $B \gets b \cdot A + c \cdot (A \cdot A)$
    \State $X \gets a \cdot X + B \cdot X$
\EndFor
\If{$\text{rows}(G) > \text{cols}(G)$}
    \State $X \gets X^\top$
\EndIf
\State \Comment{Шаг 3: момент и обновление}
\State $M \gets \beta \cdot M + (1 - \beta) \cdot X$
\State $\eta_{\text{adj}} \gets \eta \cdot 0.2 \cdot \sqrt{\max(\text{shape}(W))}$
\State $W \gets W - \eta_{\text{adj}} \cdot M - \eta \cdot \lambda \cdot W$
\State \Return $W$
\end{algorithmic}
\end{algorithm}

\subsubsection{Вычислительная сложность и память}

\begin{itemize}
    \item \textbf{Память}: $d$ скаляров на параметр (один момент $M$) + временные буферы $O(nm)$ для итераций;
    \item \textbf{Вычисления}: $O(K \cdot nm \cdot \min(n,m))$ на шаг из-за матричных умножений в итерациях Ньютона-Шульца;
    \item \textbf{Параллелизация}: матричные операции хорошо параллелятся на GPU, но итеративная природа ограничивает конвейеризацию.
\end{itemize}

\textbf{Практические замечания}:
\begin{itemize}
    \item Число итераций $K$ по умолчанию равно 5; больше итераций дают лучшее приближение, но увеличивают время;
    \item Для стабильности используется нормировка спектральной нормы и $\epsilon = 10^{-7}$;
    \item Метод не применяется к 1D-параметрам (bias, layer norm), где ортогонализация не имеет смысла.
\end{itemize}

\subsection{Hybrid: комбинированный подход}

\subsubsection{Мотивация}

Ни AdamW, ни Muon не являются универсально оптимальными:
\begin{itemize}
    \item AdamW быстр и стабилен, но может медленнее сходиться в задачах с сильными корреляциями градиентов;
    \item Muon теоретически ускоряет сходимость, но требует больше вычислений и может быть нестабилен на ранних этапах обучения.
\end{itemize}

Идея гибридного подхода: использовать сильные стороны каждого метода, назначая их разным группам параметров.

\subsubsection{Стратегия разделения параметров}

Параметры модели разделяются на две группы согласно алгоритму из \texttt{src/optimizers/optimizer.py}:

\begin{enumerate}
    \item \textbf{Фильтрация 2D-параметров}: из всех параметров отбираются только двумерные тензоры, исключая те, что содержат в имени \texttt{embed\_tokens} или \texttt{lm\_head}:
    \begin{lstlisting}[language=Python]
all_2d_params = [
    p for p in params
    if p.ndim >= 2 
    and "embed_tokens" not in str(p) 
    and "lm_head" not in str(p)
]
    \end{lstlisting}
    
    \item \textbf{Разделение пополам}: первые 50\% отобранных 2D-параметров обучаются через Muon, остальные 50\% + все 1D-параметры --- через AdamW:
    \begin{lstlisting}[language=Python]
split_idx = len(all_2d_params) // 2
muon_params = all_2d_params[:split_idx]           # первая половина → Muon
adamw_params = all_2d_params[split_idx:] + other_params  # вторая половина + остальные → AdamW
    \end{lstlisting}
\end{enumerate}

\subsubsection{Псевдокод гибридного оптимизатора}

\begin{algorithm}[H]
\caption{Hybrid: комбинированное обновление всех параметров}
\label{alg:hybrid}
\begin{algorithmic}[1]
\Require Параметры $\Theta = \{\theta_1, \ldots, \theta_N\}$, градиенты $\{g_1, \ldots, g_N\}$
\State \Comment{Фильтрация 2D-параметров}
\State $all\_2d \gets \{ p \in \Theta \mid p.\text{ndim} \ge 2 \land \text{"embed\_tokens"} \notin p \land \text{"lm\_head"} \notin p \}$
\State $other \gets \Theta \setminus all\_2d$
\State \Comment{Разделение пополам}
\State $split \gets \lfloor |all\_2d| / 2 \rfloor$
\State $muon\_params \gets all\_2d[0:split]$
\State $adamw\_params \gets all\_2d[split:] \cup other$
\State \Comment{Обновление через единый интерфейс Muon}
\State \Return \Call{Muon.step}{$muon\_params$, $adamw\_params$, $g$}
\end{algorithmic}
\end{algorithm}

\subsubsection{Преимущества и ограничения}

\textbf{Преимущества}:
\begin{itemize}
    \item \textbf{Баланс скорости и качества}: критически важные слои получают более мощную оптимизацию (Muon), остальные --- быструю (AdamW);
    \item \textbf{Гибкость}: стратегию разделения можно адаптировать под архитектуру;
    \item \textbf{Обратная совместимость}: при $split = 0$ метод вырождается в чистый AdamW.
\end{itemize}

\textbf{Ограничения}:
\begin{itemize}
    \item \textbf{Эвристика разделения}: выбор первых 50\% матриц может быть неоптимальным для всех архитектур;
    \item \textbf{Дополнительная сложность}: требуется отслеживание типов и порядков параметров;
    \item \textbf{Отсутствие теоретических гарантий}: комбинация двух оптимизаторов не имеет формального обоснования сходимости.
\end{itemize}

\section{Экспериментальная установка}
\label{sec:experiments}

\subsection{Аппаратное и программное обеспечение}

\begin{table}[H]
\centering
\caption{Аппаратная конфигурация}
\label{tab:hardware}
\begin{tabular}{lc}
\toprule
\textbf{Компонент} & \textbf{Характеристики} \\
\midrule
Процессор & Intel Core i5-11400F (6 ядер / 12 потоков) \\
Оперативная память & 16 ГБ DDR4 \\
Видеокарта & NVIDIA GeForce RTX 5060 Ti, 16 ГБ VRAM \\
Операционная система & Windows 10 Pro \\
\bottomrule
\end{tabular}
\end{table}

\begin{table}[H]
\centering
\caption{Зависимости программного обеспечения (из \texttt{requirements.txt})}
\label{tab:software}
\begin{tabular}{ll}
\toprule
\textbf{Пакет} & \textbf{Версия} \\
\midrule
torch & 2.10.0+cu128 \\
datasets & 4.8.2 \\
hydra-core & 1.3.2 \\
lm\_eval & 0.4.11 \\
matplotlib & 3.10.8 \\
mypy & 1.19.1 \\
peft & 0.18.1 \\
psutil & 7.2.2 \\
pylint & 4.0.5 \\
tokenizers & 0.22.2 \\
transformers & 5.3.0 \\
trl & 0.29.0 \\
\bottomrule
\end{tabular}
\end{table}

\noindent\textit{Примечание}: пакет \texttt{torch} установлен из дополнительного индекса \url{https://download.pytorch.org/whl/cu128} для поддержки CUDA 12.8.

\subsection{Модель и данные}

\begin{table}[H]
\centering
\caption{Параметры экспериментальной установки}
\label{tab:setup}
\begin{tabular}{ll}
\toprule
\textbf{Компонент} & \textbf{Значение} \\
\midrule
Модель & Qwen/Qwen2.5-0.5B (502M параметров) \\
Адаптация & LoRA: $r=16$, $\alpha=32$, dropout=0.05 \\
Целевые модули & q\_proj, k\_proj, v\_proj, o\_proj, gate\_proj, up\_proj, down\_proj \\
Обучаемые параметры & 8.8M (1.75\% от общего числа) \\
Тренировочный датасет & Elriggs/openwebtext-100k (1\% = 1000 сэмплов) \\
Тестовый датасет & PIQA (1838 примеров логического вывода) \\
Точность вычислений & bfloat16 (нативная поддержка на RTX 5060 Ti) \\
Размер батча & 4 на устройство, 4 шага накопления градиентов \\
Эпохи & 5 \\
Learning rate & $5 \times 10^{-5}$, cosine schedule с warmup 3\% \\
Weight decay & 0.01 \\
Оптимизаторы & AdamW / Muon / Hybrid \\
\bottomrule
\end{tabular}
\end{table}

\subsection{Метрики и протокол оценки}

Для каждого эксперимента фиксируются:
\begin{enumerate}
    \item \textbf{Время обучения}: от начала первого шага до завершения последней эпохи;
    \item \textbf{Пиковая память GPU}: максимальное значение;
    \item \textbf{Финальный loss}: среднее значение функции потерь на последнем шаге;
    \item \textbf{PIQA Accuracy}: доля правильных ответов на тестовом датасете после дообучения;
    \item \textbf{Шагов в секунду}: усреднённая скорость обработки мини-батчей.
\end{enumerate}

Все эксперименты запускаются с фиксированным random seed (42) для воспроизводимости.
\newpage

\section{Результаты}
\label{sec:results}

\subsection{Сводная таблица метрик}

\begin{table}[H]
\centering
\caption{Сравнение оптимизаторов по ключевым метрикам}
\label{tab:results}
\begin{tabular}{lccc}
\toprule
\textbf{Метрика} & \textbf{AdamW} & \textbf{Muon} & \textbf{Hybrid} \\
\midrule
Время обучения & \textbf{7м 41с} & 9м 22с & 8м 23с \\
Пиковая память, МБ & 3912 & \textbf{3846} & 3878 \\
Финальный loss & \textbf{2.863} & 2.903 & 2.888 \\
PIQA Accuracy, \% & \textbf{70.29} & 70.13 & 69.86 \\
Шагов в секунду & \textbf{0.682} & 0.560 & 0.626 \\
Trainable params & \multicolumn{3}{c}{8.8M (1.75\% от общего числа)} \\
\bottomrule
\end{tabular}
\end{table}

\subsection{Анализ производительности}

\paragraph{Скорость обучения.}
AdamW демонстрирует наилучшую производительность: 0.682 шага/с против 0.560 у Muon ($-18\%$). Это объясняется отсутствием итераций Ньютона-Шульца, требующих дополнительных матричных умножений. Гибридный подход занимает промежуточное положение (0.626 шага/с), так как лишь половина параметров обрабатывается более тяжёлым Muon.

\paragraph{Потребление памяти.}
Muon экономит $\sim$66 МБ ($\sim$1.7\%) видеопамяти по сравнению с AdamW благодаря использованию одного буфера момента вместо двух. Гибридный метод потребляет промежуточное количество памяти, пропорциональное доле параметров, обрабатываемых каждым оптимизатором.

\paragraph{Качество обучения.}
Все три метода достигают сопоставимого финального loss ($\sim$2.86--2.90) и точности на PIQA ($\sim$70\%). Небольшие различия (менее 0.5\% accuracy) статистически незначимы при данном объёме данных (1000 тренировочных сэмплов). Это указывает на то, что для задач малого масштаба выбор оптимизатора влияет в первую очередь на эффективность, а не на итоговое качество.

\paragraph{Динамика сходимости.}
Анализ логов показывает, что все оптимизаторы демонстрируют монотонное снижение loss без признаков расходимости. На ранних этапах (первые 2 эпохи) Muon иногда показывает более резкие колебания, что может быть связано с итеративной природой ортогонализации. К концу обучения траектории всех методов стабилизируются.
\newpage

\section{Заключение}
\label{sec:conclusion}

В данной работе проведён детальный сравнительный анализ трёх оптимизаторов для LoRA-дообучения больших языковых моделей:

\begin{enumerate}
    \item \textbf{AdamW} остаётся надёжным выбором для большинства практических задач: он быстр, стабилен и хорошо изучен. Его покомпонентная адаптивность эффективно справляется с разномасштабными градиентами, а коррекция weight decay улучшает обобщение.
    
    \item \textbf{Muon} предлагает инновационный подход через матричную ортогонализацию, теоретически ускоряющий сходимость в задачах с коррелированными градиентами. На практике он экономит память и показывает сопоставимое качество, но требует больше вычислений.
    
    \item \textbf{Hybrid} демонстрирует, что комбинация методов может давать сбалансированные результаты. Предложенная стратегия разделения параметров (первые 50\% 2D-матриц $\rightarrow$ Muon, остальные $\rightarrow$ AdamW) является простой, но эффективной эвристикой.
\end{enumerate}

\textbf{Ключевой вывод}: для задач малого и среднего масштаба (как в данном эксперименте) различия в итоговом качестве между оптимизаторами минимальны. Выбор следует делать исходя из ограничений по времени и памяти: AdamW для скорости, Muon для экономии ресурсов, Hybrid для баланса.

\textbf{Практическая рекомендация}: начинайте с AdamW как с базового варианта. При наличии избыточного времени или требований к минимизации памяти экспериментируйте с Muon или гибридными схемами, обязательно проводя валидацию на отложенной выборке.

\section*{Список литературы}
\label{sec:references}

\begin{thebibliography}{9}

\bibitem{loshchilov2019adamw}
Loshchilov I., Hutter F. 
\textit{Decoupled Weight Decay Regularization}. 
International Conference on Learning Representations (ICLR), 2019.

\bibitem{hu2021lora}
Hu E. J. et al. 
\textit{LoRA: Low-Rank Adaptation of Large Language Models}. 
arXiv preprint arXiv:2106.09685, 2021.

\bibitem{moonshot2024muon}
MoonshotAI. 
\textit{Muon: Matrix Orthogonalization for Efficient LLM Training}. 
GitHub repository, 2024. 
\url{https://github.com/MoonshotAI/Moonlight}

\bibitem{bisk2020piqa}
Bisk Y. et al. 
\textit{PIQA: Reasoning about Physical Commonsense in Natural Language}. 
Proceedings of the AAAI Conference on Artificial Intelligence, 2020.

\bibitem{mezo2023}
Malladi S. et al. 
\textit{MeZO: Memory-Efficient Zeroth-Order Optimizer for LLM Fine-tuning}. 
arXiv preprint arXiv:2305.17333, 2023.

\bibitem{schulz1933newton}
Schulz G. 
\textit{Iterative Berechnung der reziproken Matrix}. 
Zeitschrift für Angewandte Mathematik und Mechanik, 1933.

\end{thebibliography}

\appendix
\section*{Приложение: Исходные логи экспериментов}

Полные логи обучения и оценки доступны в репозитории:
\begin{itemize}
    \item \texttt{outputs/checkpoints/metrics\_adamw.json}
    \item \texttt{outputs/checkpoints/metrics\_muon.json}
    \item \texttt{outputs/checkpoints/metrics\_hybrid.json}
    \item \texttt{logs/train\_*.log}
\end{itemize}

Для воспроизведения результатов выполните:
\begin{lstlisting}[language=bash]
# Клонирование репозитория
git clone https://github.com/Not-broken-today/llm-training-acceleration
cd llm-training-acceleration

# Установка зависимостей
pip install -r requirements.txt --extra-index-url https://download.pytorch.org/whl/cu128

# Запуск эксперимента (пример для AdamW)
bash run_file/AdamW_run.sh

# Проверка кода
bash pylint_test_code.sh
mypy src/
\end{lstlisting}

\end{document}